% =============================================================================
\chapter{Contribuições Esperadas}
% =============================================================================

\begin{itemize}
\item uma arquitetura formal integrando HTN, fontes dinâmicas de preferências e MORL;
\item um método para empregar MORL especificamente na seleção de métodos de decomposição HTN;
\item uma separação explícita entre restrições simbólicas obrigatórias e preferências comportamentais flexíveis;
\item uma política condicionada a ser avaliada quanto à variação de comportamento sem retreinamento por vetor de pesos;
\item um mecanismo rastreável de decisão hierárquica para NPCs adaptativos;
\item uma comparação experimental entre seleção aprendida e heurísticas tradicionais de HTN;
\item uma comparação controlada entre perfis, regras, grafos relacionais e,
quando viável, codificadores profundos de preferências;
\item evidências sobre benefícios, limitações e custo da integração neuro-simbólica proposta;
\item um framework experimental reutilizável para pesquisas futuras em Game AI híbrida.
\end{itemize}

% =============================================================================
\chapter{Considerações Finais}
% =============================================================================

A reformulação apresentada resolve a principal ambiguidade da hipótese
anterior.
O MORL não é responsável por execução motora nem pela escolha direta de ações
primitivas.
Ele atua no ponto em que um planejador HTN possui mais de um método semanticamente
aplicável e precisa escolher qual decomposição melhor atende às prioridades correntes.

A arquitetura pode ser resumida da seguinte forma: o HTN determina o que é
válido; uma fonte de preferências fornece $\mathbf{w}_t$ para expressar o
compromisso corrente; e o MORL seleciona o método válido mais apropriado.
A geração de preferências é um componente auxiliar e comparável: perfis e
regras são baselines; um grafo relacional inspirado em GOAP e um modelo profundo
são implementações opcionais da fonte, condicionadas ao orçamento e aos resultados.
A execução das tarefas primitivas permanece sob responsabilidade de sistemas
convencionais.

O baseline HTN existente é compatível com esse direcionamento, mas deverá ser
adaptado para expor o conjunto de métodos aplicáveis e aceitar um seletor
externo.
A avaliação deverá demonstrar não apenas ganho de desempenho, mas também
generalização entre preferências, validade semântica, qualidade das explicações
e custo compatível com decisões em tempo real.

\vspace{3em}
\noindent Porto Alegre, agosto de 2026.

\vspace{4em}
\begin{flushleft}
\rule{7cm}{0.4pt}\\
\textbf{Anderson G. A. C. Filho}\\
Mestrando em Ciência da Computação --- PPGC/UFRGS
\end{flushleft}
